


\noindent\textbf{Source and scope.}
The source read is Alain Connes and Caterina Consani, \emph{The Absolute
Twistor Line and the Geometry of $\overline{\operatorname{Spec}\Z}$},
arXiv:2609.00299v1, original author source \nolinkurl{CC.tex}, 1228 lines,
SHA--256
\nolinkurl{57a10dcef5cd758e6b2f1c54b1c0a11b0fb093638f74fb1be515ec3bc7080ba4}.
Reading coverage is the complete file, including its bibliography.
The signed base, imaginary chart, geometric involution, and odd maps are
Connes--Consani's constructions, not new programme constructions.
All assertions below are proved in the specified category of pointed
commutative monoids, or in ordinary algebraic $K_0(\PP^1_\C)$.
No claim about a more general category of spherical algebras is inferred
merely from this monoid calculation.

\section{The actual base and tensor product}

Let
\begin{equation}\tag{HA1}
 B=\{0,1,\eps\},\qquad \eps^2=1,\quad \eps\ne1,
 \qquad D=B[J]/(J^2=\eps).
\end{equation}
Here $0$ is absorbing and each displayed presentation is a presentation
of a pointed commutative monoid. In particular there is no additive
relation $1+\eps=0$. A $B$-monoid is a pointed commutative monoid with a
specified map from $B$, and a $B$-morphism fixes that map. Put
\begin{equation}\tag{HA2}
 C_+=B[T],\quad C_\times=B[T,T^{-1}],\qquad
 A_+=D[T],\quad A_\times=D[T,T^{-1}].
\end{equation}
These are the spatial chart and overlap, and their imaginary enrichments.
The relation in (HA1) is the one in \CCsource{576--600}.

The tensor product $P\otimes_B Q$ means the smash product of pointed
monoids modulo $(bp)\otimes q=p\otimes(bq)$ for $b\in B$.
Its multiplication is
$(p\otimes q)(p'\otimes q')=pp'\otimes qq'$. In particular
\begin{equation}\tag{HA3}
 \eps\otimes1=1\otimes\eps=\eps,
 \qquad \eps\otimes\eps=1.
\end{equation}
These are exactly the balancing relations specified in
\CCsource{437--465}.

\begin{lemma}
Every nonzero element of $A_+$ has a unique expression $J^aT^b$ with
$a\in\Z/4\Z$ and $b\in\N$. For $A_\times$, replace $\N$ by $\Z$.
The units of $A_+\otimes_B A_+$ have the presentation
\[
 \langle j_1,j_2,\eps\mid
 j_1^2=j_2^2=\eps,\ \eps^2=1,\ [j_1,j_2]=1\rangle
 \cong (\Z/4\Z)\times(\Z/2\Z).
\]
In this tensor product $\eps\ne1$.
\end{lemma}
\begin{proof}
The pointed monoid obtained by adjoining an absorbing zero to
$(\Z/4\Z)\times\N$, with $J=(1,0)$, $T=(0,1)$ and
$\eps=(2,0)$, satisfies (HA1)--(HA2). Conversely those presentations
reduce uniquely to these coordinates. The same construction works with
$\Z$ in the second coordinate. In the tensor product use
$j_1=J\otimes1$, $j_2=1\otimes J$. The only additional relation between
the two factors is $j_1^2=j_2^2$. Setting $h=j_2j_1^{-1}$ gives
$h^2=1$; the inverse change of generators is $j_2=j_1h$.
Thus the unit group is the asserted direct product. Its element
$\eps=j_1^2$ has order two. Alternatively, the compatible pair of maps
to the multiplicative monoid of $\C$, sending both $J$'s to $i$ and
$\eps$ to $-1$, induces a map on the balanced tensor product which
separates $\eps$ and $1$.
\end{proof}

\begin{proposition}[The imaginary diagonal does not descend]
The assignments
\begin{equation}\tag{HA4}
 \Delta(T)=T\otimes T,\qquad
 \Delta(J)=J\otimes J,\qquad \Delta(\eps)=\eps
\end{equation}
do not define a $B$-morphism on $A_+$ or $A_\times$.
Moreover, neither of those two $B$-monoids admits a counit to $B$.
\end{proposition}
\begin{proof}
A $B$-morphism must preserve $J^2=\eps$. But (HA3) gives
\begin{equation}\tag{HA5}
 (J\otimes J)^2=\eps\otimes\eps=1\ne\eps
       =\Delta(J^2).
\end{equation}
For the second assertion, the image of the unit $J$ under a unital monoid
map is a unit. The only units of $B$ are $1$ and $\eps$, both with square
$1$. Neither squares to $\eps$. Consequently there is no $B$-morphism
from either enriched chart to $B$ at all. A counital comonoid structure,
and hence a Hopf structure over $B$, would include such a morphism.
\end{proof}

The scope matters: the coproduct in the author's definition is on the
\emph{spatial} monoid $M$ of $B[M]$ (\CCsource{462--465}), not an
assertion that a diagonal on an additional generator survives every
signed quotient relation. Equation (HA5) identifies the exact relation
which prevents that extension; it does not erase the spatial coproduct.

\section{The surviving coaction and the stronger torsor correspondence}

\begin{proposition}[Spatial coaction]
There is a counital and coassociative $B$-morphism
\begin{equation}\tag{HA6}
 \rho:A_+\longrightarrow C_+\otimes_B A_+,
 \qquad T\longmapsto T\otimes T,\quad J\longmapsto1\otimes J.
\end{equation}
It extends to $A_\times\to C_\times\otimes_B A_\times$.
The spatial bialgebra has $\Delta_C(T)=T\otimes T$ and counit
$\eta_C(T)=1$. On $C_\times$ the antipode is $T\mapsto T^{-1}$.
\end{proposition}
\begin{proof}
The image of $J$ in (HA6) has square $1\otimes\eps=\eps$, as required.
There is no other relation to check beyond invertibility of $T$ on the
overlap, where its image is invertible. Both coassociativity composites
send $T$ to $T\otimes T\otimes T$ and $J$ to $1\otimes1\otimes J$.
Applying the spatial counit sends them back to $T$ and $J$.
The assertions about $C$ follow by the same checks on $T$ and the base.
\end{proof}

For a commutative $B$-monoid $R$ let
\[
 X(R)=\Hom_B(A_\times,R)
     =\{(t,j)\in R^\times\times R^\times:j^2=\eps_R\}.
\]
The spatial coaction is the actual map
\begin{equation}\tag{HA7}
 R^\times\times X(R)\longrightarrow X(R),
 \qquad (u,(t,j))\longmapsto(ut,j).
\end{equation}
This includes, on complex points, $\C^\times$ acting without changing the
chosen root of $-1$. It proves the precise action on the unfolded chart.

There is a stronger intrinsic relation. Define
\begin{equation}\tag{HA8}
 H=B[u,u^{-1},h]/(h^2=1),\qquad
 G(R)=R^\times\times\{h\in R^\times:h^2=1\}.
\end{equation}
The letter $h$ in (HA8) is an independent generator; it is not the base
element $\eps$. The Hopf structure on $H$ sends $u$ and $h$ to their
tensor diagonals, with counit $u,h\mapsto1$, and antipode
$u\mapsto u^{-1}$, $h\mapsto h$.

\begin{theorem}[Exact torsor identities and ternary product]
There is a coaction and a canonical isomorphism
\begin{align}
 \rho_H:A_\times&\longrightarrow H\otimes_B A_\times,
 &T&\longmapsto u\otimes T,&J&\longmapsto h\otimes J,
 \tag{HA9}\\
 \operatorname{can}:A_\times\otimes_B A_\times
   &\xrightarrow{\ \sim\ } H\otimes_B A_\times,
 &&&&\tag{HA10}
\end{align}
where $\operatorname{can}$ sends
\[
 T_1\mapsto uT,\quad T_2\mapsto T,
 \qquad J_1\mapsto hJ,\quad J_2\mapsto J.
\]
When $X(R)$ is nonempty, the action
$(u,h)(t,j)=(ut,hj)$ makes it a principal homogeneous set for $G(R)$.
The ternary operation below requires no choice of a root:
\begin{equation}\tag{HA11}
 [x,y,z]=(t_x t_y^{-1}t_z,\ j_x j_y^{-1}j_z),
 \qquad X(R)^3\longrightarrow X(R).
\end{equation}
This operation satisfies $[x,y,y]=x=[y,y,x]$ and
$[[x,y,z],u,v]=[x,y,[z,u,v]]$.
\end{theorem}
\begin{proof}
For (HA9), $(h\otimes J)^2=1\otimes\eps=\eps$;
coassociativity and the counit follow on each generator.
The inverse to (HA10) is the explicitly defined $B$-morphism
\begin{equation}\tag{HA12}
 u\longmapsto T_1T_2^{-1},\quad
 h\longmapsto J_1J_2^{-1},\quad
 T\longmapsto T_2,\quad J\longmapsto J_2.
\end{equation}
Indeed $(J_1J_2^{-1})^2=\eps\eps^{-1}=1$, so (HA12)
respects the defining equation of $h$; each composite fixes the displayed
generators. Given $(t,j),(t',j')\in X(R)$, the unique group element
carrying the second to the first is
$(tt'^{-1},jj'^{-1})$, whose second coordinate has square $1$.
The square of the second coordinate in (HA11) is
$\eps_R\eps_R^{-1}\eps_R=\eps_R$.
The stated identities follow by cancellation in the group of units of
$R$. In particular they are identities of the actual functor $X$, not
only of its complex points. The coordinate morphism for (HA11) is
$T\mapsto T_1T_2^{-1}T_3$, $J\mapsto J_1J_2^{-1}J_3$.
\end{proof}

Here ``torsor identities'' means the proved coaction and canonical-map
isomorphism (HA9)--(HA12); no unspecified Grothendieck topology or
descent theorem is being used.

\begin{proposition}[The precise effect of choosing an imaginary root]
After the explicit base extension $B\to D$, let $j_0$ denote the root
in the base copy of $D$. There is an isomorphism of $D$-monoids
\begin{equation}\tag{HA13}
 D\otimes_B A_\times
       \cong D[T,T^{-1},k]/(k^2=1),
 \qquad k=Jj_0^{-1}.
\end{equation}
Transporting the usual Hopf structure on the right gives
\begin{equation}\tag{HA14}
 \Delta_D(J)=j_0^{-1}(J\otimes J),\quad
 \eta_D(J)=j_0,\quad \mathcal S_D(J)=J,
 \qquad \Delta_D(T)=T\otimes T.
\end{equation}
Replacing $j_0$ by $\eps j_0$ changes the displayed coproduct on $J$
by the factor $\eps$ and changes its counit by the same factor.
\end{proposition}
\begin{proof}
One has $k^2=J^2j_0^{-2}=\eps\eps^{-1}=1$.
The inverse to (HA13) sends $J$ to $j_0k$.
The Hopf formula $\Delta(k)=k\otimes k$ gives
$\Delta(J)=j_0 k\otimes k=j_0^{-1}J\otimes J$;
the counit gives $\eta(J)=j_0$. Since $k^{-1}=k$, the antipode
fixes $J=j_0k$. The last assertion follows from $\eps^{-1}=\eps$.
This calculation also shows explicitly what root choice the missing
$B$-counit would have had to make.
\end{proof}

\section{The quotient symmetry and the spatial action}

On the unfolded complex overlap the source defines
\begin{equation}\tag{HA15}
 \alpha(z,j)=(-1/z,-j),\qquad j\in\{i,-i\},
\end{equation}
by $\alpha^*(T)=\eps T^{-1}$, $\alpha^*(J)=\eps J$
(\CCsource{773--806}). If $r_\lambda(z,j)=(\lambda z,j)$,
then
\begin{equation}\tag{HA16}
 \alpha r_\lambda=r_{\lambda^{-1}}\alpha.
\end{equation}
This is verified by evaluating both sides at $(z,j)$: both give
$(-\lambda^{-1}/z,-j)$. Thus the spatial action and involution have an
exact semidirect-product relation.

The same uniform map $r_\lambda$ does not descend to the orbit quotient
for arbitrary $\lambda$. Indeed two representatives $(z,j)$ and
$(-1/z,-j)$ are sent to $(\lambda z,j)$ and $(-\lambda/z,-j)$.
The latter is in the orbit of the former, for all $z\ne0,\infty$,
exactly when $-\lambda/z=-1/(\lambda z)$, that is, $\lambda^2=1$.
After choosing the representative component $j=i$, the action which
\emph{does} lift the usual quotient scaling is
\begin{equation}\tag{HA17}
 \widetilde r_\lambda(z,i)=(\lambda z,i),\qquad
 \widetilde r_\lambda(z,-i)=(\lambda^{-1}z,-i).
\end{equation}
Its commutation with $\alpha$ follows from (HA16), and its multiplication
law follows component by component. This supplies the exact additional
branch rule needed to transport scaling through the geometric fold.
It does not replace the underlying spatial-core coaction by a diagonal
on $J$.

\section{Odd Frobenius maps on the projective line}

For a positive odd integer $n$, define
\begin{equation}\tag{HA18}
 \chi(n)=(-1)^{(n-1)/2},\quad e(n)=\frac{1-\chi(n)}2,
 \quad a(z)=-\frac1z,\quad p_n(z)=z^n,
 \quad F_n=a^{e(n)}p_n.
\end{equation}
Thus $F_n(z)=z^n$ for $n\equiv1\pmod4$ and
$F_n(z)=-1/z^n$ for $n\equiv3\pmod4$, as in Connes--Consani's
Proposition \texttt{prop:twistor\_frobenius}
(\CCsource{1159--1191}).
In homogeneous coordinates $z=X/Y$, these are everywhere-defined maps
\begin{equation}\tag{HA19}
 F_n[X:Y]=
 \begin{cases}
 [X^n:Y^n],&\chi(n)=1,\\
 [-Y^n:X^n],&\chi(n)=-1.
 \end{cases}
\end{equation}
The two coordinate polynomials never vanish simultaneously on
$\PP^1$, so each map has degree $n$.

\begin{theorem}[Composition, real structure, and the exact equivariance]
For all positive odd $m,n$,
\begin{equation}\tag{HA20}
 F_mF_n=F_{mn},\qquad F_1=\id.
\end{equation}
For $\sigma(z)=-1/\overline z$ one has
\begin{equation}\tag{HA21}
 F_n\sigma=\sigma F_n.
\end{equation}
For the usual scaling action of $\C^\times$ the exact formula is
\begin{equation}\tag{HA22}
 F_n(\lambda z)=\theta_n(\lambda)F_n(z),\qquad
 \theta_n(\lambda)=\lambda^{n\chi(n)}.
\end{equation}
\end{theorem}
\begin{proof}
For odd $n$, $p_na(z)=(-1/z)^n=-1/z^n=ap_n(z)$ on the
nonzero affine coordinate, hence on $\PP^1$ by (HA19).
Also $a^2=\id$ and $p_mp_n=p_{mn}$. The identity
$\chi(mn)=\chi(m)\chi(n)$ follows by checking the two invertible
residue classes modulo four. Consequently
$e(mn)\equiv e(m)+e(n)\pmod2$, proving (HA20).
Complex conjugation $c$ commutes with $a$ and each $p_n$, which have
real coefficients; $\sigma=ac$ gives (HA21).
Finally (HA22) follows by substitution in the two cases of (HA18).
The homogeneous formulas prove all these equalities at $0$ and $\infty$
as well.
\end{proof}

\section{A complete computation of the ordinary \texorpdfstring{$K_0$}{K0} action}

Here $K_0(\PP^1_\C)$ is the Grothendieck ring of algebraic vector bundles,
with addition from exact sequences and multiplication from tensor product.
Let $L=[\OO(1)]$ and $\eta=L-1$.

\begin{lemma}
\begin{equation}\tag{HA23}
 K_0(\PP^1_\C)=\Z[\eta]/(\eta^2),\qquad
 [\OO(d)]=1+d\eta\quad(d\in\Z).
\end{equation}
Every vector bundle $E$ has
\begin{equation}\tag{HA24}
 [E]=r+q\eta,\qquad r=\rank E,\qquad
 q=\chi(E)-r.
\end{equation}
\end{lemma}
\begin{proof}
The Euler sequence on $\PP^1$ is
$0\to\OO(-1)\to\OO^{\oplus2}\to\OO(1)\to0$;
its maps in homogeneous coordinates are $(-Y,X)$ and $(X,Y)$.
The composition is zero, and checking the charts $X\ne0$ and $Y\ne0$
shows that its kernel and cokernel have the asserted form. Hence
$L+L^{-1}=2$ and $(L-1)^2=0$.

For completeness, generation of $K_0$ by $1,L^{-1}$ follows without
assuming that every bundle has already been split. On
$\PP^1\times\PP^1$, with projections $q_1,q_2$, the diagonal is cut out
by $X_1Y_2-Y_1X_2$, giving the exact sequence
\[
 0\longrightarrow\OO(-1,-1)\longrightarrow\OO
   \longrightarrow\OO_\Delta\longrightarrow0.
\]
Tensor it by $q_1^*E$ and apply the alternating direct image under $q_2$.
The diagonal term is $E$. The other two terms have direct images
$H^j(E)\otimes\OO$ and $H^j(E(-1))\otimes\OO(-1)$ respectively.
These identities can be checked by the two-chart Cech complex on the
first factor, tensored with either indicated line bundle on the second;
there are no terms in degrees greater than one. The alternating
identity of the resulting long exact sequence gives
\begin{equation}\tag{HA25}
 [E]=\chi(E)[\OO]-\chi(E(-1))[\OO(-1)].
\end{equation}
Thus $1,L^{-1}$ generate $K_0$, and so do $1,\eta$.
The same two-chart Cech calculation gives
$H^0(\OO)=\C$, $H^1(\OO)=0$, and
$H^0(\OO(-1))=H^1(\OO(-1))=0$.
Rank and Euler characteristic therefore map $[\OO],[\OO(-1)]$ to
$(1,1),(1,0)$ in $\Z^2$. Their determinant is $-1$, so there are
no additional additive relations. This proves the ring presentation.
For $d\ge0$, $(1+\eta)^d=1+d\eta$; for negative $d$ use
$(1+\eta)^{-1}=1-\eta$. Finally rank sends $1,\eta$ to $1,0$,
while Euler characteristic sends them to $1,1$ (the Euler sequence
gives $\chi(\OO(1))=2$). Formula (HA24) follows.
\end{proof}

\begin{theorem}[Exact ordinary Adams equality]
For every positive odd $n$,
\begin{equation}\tag{HA26}
 F_n^*(1+q\eta)=1+nq\eta,\qquad
 F_n^*(r+q\eta)=r+nq\eta,
 \qquad F_n^*=\psi^n\quad\hbox{on }K_0(\PP^1_\C).
\end{equation}
\end{theorem}
\begin{proof}
The two homogeneous coordinates of (HA19) are sections of $\OO(n)$
with no common zero. The universal tautological quotient defining the
map therefore gives $F_n^*\OO(1)\simeq\OO(n)$. This can equally be
checked by pulling back the transition ratio $X/Y$ on the inverse images
of the two standard charts, with their interchange in the second case
of (HA19). Thus $F_n^*\eta=L^n-1=n\eta$, proving the two explicit
formulas.

Define the Adams operations from the usual exterior-power operations
by the formal identity
\[
 \frac{d}{dt}\log\lambda_t(x)
   =\sum_{n\ge1}(-1)^{n-1}\psi^n(x)t^{n-1},
 \qquad \lambda_t(x+y)=\lambda_t(x)\lambda_t(y).
\]
It shows additivity of each $\psi^n$. For a line-bundle class $M$,
$\lambda_t(M)=1+Mt$, so $\psi^n(M)=M^n$.
The additive generators $1,L^{-1}$ from (HA25) therefore have identical
images under $F_n^*$ and $\psi^n$, proving the equality on every class.
In particular the holomorphic involution $a$ acts trivially on this
$K_0$, because $a^*\OO(1)\simeq\OO(1)$.
\end{proof}

This is the exact ordinary-$K_0$ statement justified by the maps. It
proves a comparison on the receiving group in (HA23), rather than an
identity of all additional structures which a bundle may carry.

\section{What happens to equivariant weights and real structures}

Let $E$ be a $\C^\times$-equivariant bundle for the usual scaling action
on $\PP^1$. Pullback by $F_n$ acquires a specified equivariant structure
using the group homomorphism $\theta_n$ in (HA22):
\begin{equation}\tag{HA27}
 \lambda\cdot(z,v)=(\lambda z,\theta_n(\lambda)\cdot v),
 \qquad v\in E_{F_n(z)}.
\end{equation}
Indeed the right-hand side belongs to $E_{F_n(\lambda z)}$ by (HA22),
and associativity follows from multiplicativity of $\theta_n$.
If the character weights at $0$ and $\infty$ are the multisets
$W_0,W_\infty$, the exact transformed multisets are
\begin{equation}\tag{HA28}
 (W'_0,W'_\infty)=
 \begin{cases}
 (nW_0,nW_\infty),&n\equiv1\pmod4,\\
 (-nW_\infty,-nW_0),&n\equiv3\pmod4.
 \end{cases}
\end{equation}
The fibre identification is
$(F_n^*E)_z=E_{F_n(z)}$; $F_n$ fixes the poles in the first case and
interchanges them in the second. A target character $\lambda^m$
becomes $(\lambda^{n\chi(n)})^m$, which proves (HA28).

There is an explicit test showing that (HA26) cannot simply be promoted
to the same assertion about this equivariant pullback. Give the trivial
line bundle the fibre action of character $\lambda$. Under (HA27) its
character becomes $\lambda^{n\chi(n)}$. Its equivariant Adams image
has character $\lambda^n$. For $n\equiv3\pmod4$ these are different:
restricting to the fixed point $0$ detects the distinct Laurent monomials
$t^{-n}$ and $t^n$ in the representation ring
$R(\C^\times)=\Z[t,t^{-1}]$. Both forget to the identical class $1$
in the ordinary $K_0$ used in (HA26).

Suppose now that a bundle is supplied with an actual antilinear lift
$\tau:E_z\to E_{\sigma(z)}$ of the antipodal real structure. The
pullback has the exact antilinear lift
\begin{equation}\tag{HA29}
 \tau_n(z,v)=(\sigma(z),\tau v).
\end{equation}
Its target fibre is correct because
$\sigma F_n(z)=F_n\sigma(z)$ by (HA21).
If the supplied lift squares to $+\id$ or $-\id$, its pullback squares
to the same operator. This is a construction from an existing lift;
it does not conjure one from the class of the underlying vector bundle.
For example an actual isomorphism $E\simeq\OO(w)^{\oplus r}$ pulls back
to $\OO(nw)^{\oplus r}$, while the rank and degree calculation alone
does not specify a filtration, a polarization, a connection, or the
equivariant fibre weights. Equations (HA27)--(HA29) specify precisely
the further structures transported here.

The source's discussion of real Hodge structures and Adams operations
appears in \CCsource{416--417} and in its section \emph{Frobenius and
Twistor Equivariance}. The calculation here supplies the exact
ordinary-$K_0$ map, the character correction for the second residue
class, and the transport of any supplied real lift. No equality with a
programme cohomology group, no positivity theorem for a trace pairing,
and no Riemann-hypothesis conclusion is asserted by these identities.

\begin{thebibliography}{9}
\bibitem{HA-CC}
Alain Connes and Caterina Consani,
\emph{The Absolute Twistor Line and the Geometry of
$\overline{\operatorname{Spec}\Z}$},
arXiv:2609.00299v1 (2026), original author LaTeX source
\nolinkurl{CC.tex}.
\href{https://arxiv.org/abs/2609.00299v1}{Versioned paper identity};
\href{https://arxiv.org/src/2609.00299v1}{original source archive}.
The exact author-source line locations used above refer to the SHA--256
edition identified at the beginning of this derivation.
\end{thebibliography}
